\documentclass[manuscript,screen,review,sigplan]{acmart}
\settopmatter{printfolios=true}

%%
%% \BibTeX command to typeset BibTeX logo in the docs
\AtBeginDocument{%
  \providecommand\BibTeX{{%
    Bib\TeX}}}

%% Rights management information.  This information is sent to you
%% when you complete the rights form.  These commands have SAMPLE
%% values in them; it is your responsibility as an author to replace
%% the commands and values with those provided to you when you
%% complete the rights form.
\setcopyright{acmlicensed}
\copyrightyear{2026}
\acmYear{2026}
\acmDOI{XXXXXXX.XXXXXXX}
%% These commands are for a PROCEEDINGS abstract or paper.
\acmConference[MPLR2026]{The 23rd International Conference on Managed Programming Languages \& Runtimes}{Mon 29 June - Fri 3 July 2026}{Brussels, Belgium}
%%
%%  Uncomment \acmBooktitle if the title of the proceedings is different
%%  from ``Proceedings of ...''!
%%
%%\acmBooktitle{Woodstock '18: ACM Symposium on Neural Gaze Detection,
%%  June 03--05, 2018, Woodstock, NY}
\acmISBN{XXXXXXX}


\usepackage{listings}
%% auto break lines
\lstset{breaklines=true}

\usepackage{xspace}
\usepackage{pgfplots}
\usepackage[endianness=big]{bytefield}
\usetikzlibrary{patterns}
\pgfplotsset{compat=1.18}

\newcommand{\bitfieldwidth}{0.045\columnwidth}

\newcommand{\code}[1]{\colorbox{gray!10}{\texttt{$#1$}}}
\newcommand{\st}[1]{\code{#1}}
\newcommand{\js}{JavaScript\xspace}
\newcommand{\IV}{\emph{InstVar}\xspace}
\newcommand{\IVs}{\emph{InstVars}\xspace}
\input{aarch64}
\input{x86}
\input{pareto}
\newcommand{\benchmarkBarChart}[3]{%
\begin{tikzpicture}
\begin{axis}[
    ybar,
    bar width=5pt,
    width=\columnwidth,
    height=\columnwidth,
    ylabel={Execution Time (ms)},
%    xlabel={Benchmark},
    symbolic x coords={IntegerBr, Float},
    xtick=data,
    ymin=0,
    ymax=4900, %#1,
    legend style={
        at={(0.02,0.98)},
        anchor=north west,
        legend columns=5,
        font=\tiny,
        draw=black,
    },
    grid=major,
    grid style={line width=0.2pt, draw=gray!30},
    ymajorgrids=true,
    xmajorgrids=false,
    enlarge x limits=0.4,
]
#3
\legend{#2}
\end{axis}
\end{tikzpicture}
}

\begin{document}
\section{Experimental Design and Results}\label{results}
\subsection{Results}
\iftrue
Table~\ref{tab:aarch64T} shows the total execution time of the various encodings to execute the benchmark  (\code{40~fibonacci}) for integer and floating-point values on the Apple M1.

\begin{table}[htbp]
    \centering
    \aarchTotal
    \caption{AArch64 (Apple M1) total timing for encodings}
    \label{tab:aarch64T}
\end{table}

\begin{figure}[htbp]
    \centering
    \aarchBarChart
    \caption{AArch64 benchmark results}
    \label{fig:aarch64}
\end{figure}
\begin{figure}[htbp]
    \centering
    \xBarChart
    \caption{x86\_64 benchmark results}
    \label{fig:x86}
\end{figure}

Figures \ref{fig:aarch64} and \ref{fig:x86} show at a glance the relative performance of the various encodings on the two architectures.

  All timings are the median of 10 runs following 3 runs of warmups.
They are run on Linux systems running in console mode.
We calculate the mean, median and geometric mean for each example and they are essentially equal in all cases.

\begin{table*}[htbp]
      \centering
      \begin{minipage}{0.48\textwidth}
          \centering
          \aarchSummary
          \caption{AArch64 (Apple M1) net timing for encodings}
          \label{tab:aarch64}
      \end{minipage}
      \hfill
      \begin{minipage}{0.48\textwidth}
          \centering
          \xSummary
          \caption{x86-64 (Linux) net timing for encodings}
          \label{tab:x86}
      \end{minipage}
\end{table*}

Tables~\ref{tab:aarch64} and ~\ref{tab:x86} present the net execution times for each encoding across the two benchmark configurations on AArch64 (Apple M1) and x86-64 respectively. Each row reports the time for a single encoding under two conditions: \textbf{IntegerBr} (inlined branching integer Fibonacci), and \textbf{Float} (inlined float Fibonacci).
\fi
\subsubsection{Pareto Analysis}
The Pareto frontier contains the points that represent optimal trade-offs between two parameters.
In this case, we are looking at integer versus floating-point performance in figure~\ref{pareto}.
The Pareto frontiers don't intersect, showing that the Aarch64 dominates the x86-64 in this experiment.
\begin{figure}[htbp]
    \centering
    \aarchScatterPlot
    \caption{IntegerBr vs Float - AARCH64}
    \label{paretoA}
\end{figure}
\begin{figure}[htbp]
    \centering
    \xScatterPlot
    \caption{IntegerBr vs Float - X86-64}
    \label{paretoX}
\end{figure}
\begin{figure}[htbp]
    \centering
    \scatterPlotRaw
    \caption{IntegerBr vs Float}
    \label{pareto}
\end{figure}
\iffalse
\newpage
\appendix
\section{Appendix - Raw samples}
\subsection*{AARCH64 raw execution}
\aarchAppendix
\subsection*{X86-64 raw execution}
\xAppendix
\fi
\end{document}
\endinput
